
\documentclass[11pt,a4paper]{article}

\usepackage[a4paper,top=0.78in,bottom=0.72in,left=0.78in,right=0.78in]{geometry}
\usepackage[T1]{fontenc}
\usepackage[utf8]{inputenc}
\usepackage{lmodern}
\usepackage{microtype}
\usepackage{setspace}
\usepackage{graphicx}
\usepackage{float}
\usepackage{booktabs}
\usepackage{array}
\usepackage{longtable}
\usepackage{tabularx}
\usepackage{enumitem}
\usepackage{fancyhdr}
\usepackage{titlesec}
\usepackage{tocloft}
\usepackage[colorlinks=true,linkcolor=black,urlcolor=blue,citecolor=black]{hyperref}
\usepackage{xcolor}
\usepackage{xurl}
\usepackage{caption}
\usepackage{parskip}
\usepackage{lastpage}

\setstretch{1.08}
\setlength{\parindent}{0pt}
\setlength{\parskip}{5pt}
\setlist[itemize]{leftmargin=1.5em,itemsep=2pt,topsep=2pt}
\setlist[enumerate]{leftmargin=1.7em,itemsep=2pt,topsep=2pt}

\definecolor{darkblue}{RGB}{25,55,85}
\definecolor{lightgray}{RGB}{245,246,248}

\pagestyle{fancy}
\fancyhf{}
\lhead{\small\textit{InternNova – Week 4: Vulnerability Assessment \& Penetration Testing}}
\renewcommand{\headrulewidth}{0.4pt}
\cfoot{\thepage}

\titleformat{\section}{\Large\bfseries}{\thesection}{0.6em}{}
\titleformat{\subsection}{\large\bfseries}{\thesubsection}{0.6em}{}
\titleformat{\subsubsection}{\normalsize\bfseries}{\thesubsubsection}{0.6em}{}

\captionsetup[figure]{font=small,labelfont=bf,justification=centering}
\captionsetup[table]{font=small,labelfont=bf,justification=centering}

\newcommand{\refbasis}[2]{\textbf{Reference Basis:} #1\\\url{#2}\par}
\newcommand{\refstd}[2]{\textbf{Reference / Standard:} #1\\\url{#2}\par}
\newcommand{\video}[2]{\textbf{Learning Video:} #1\\\url{#2}\par}
\newcommand{\fig}[3][0.78\textwidth]{%
\begin{figure}[H]
\centering
\includegraphics[width=#1,height=0.42\textheight,keepaspectratio]{images/#2}
\caption{#3}
\end{figure}}

\begin{document}

% ---------------- COVER ----------------
\begin{titlepage}
\thispagestyle{empty}
\centering
\vspace*{1.15cm}
{\Huge\bfseries InternNova Internship Program\par}
\vspace{0.35cm}
{\Large\bfseries Week 4 Assignment\par}
\vspace{0.35cm}
{\Large Cyber Security – 6 Week Batch\par}
\vspace{1.2cm}

\rule{0.82\textwidth}{0.7pt}\\[0.35cm]
{\LARGE\bfseries Vulnerability Assessment \&\\[0.15cm] Penetration Testing\par}
\vspace{0.35cm}
\rule{0.82\textwidth}{0.7pt}

\vfill
\begin{tabular}{>{\bfseries}p{4.2cm}p{6.5cm}}
Intern Name & NAJMUDEEN S\\
Internship & InternNova\\
Week & 4\\
Assignment & Vulnerability Assessment \& Penetration Testing\\
Month & October 2026\\
\end{tabular}
\vfill

{\large\bfseries Cyber Security Internship Report\par}
\end{titlepage}

% ---------------- TOC ----------------
\pagenumbering{roman}
\setcounter{page}{2}
\tableofcontents
\newpage

% ---------------- MAIN ----------------
\pagenumbering{arabic}
\setcounter{page}{1}

\section{Task 1: Understanding Vulnerabilities}

\subsection{Introduction}
A vulnerability is a weakness in a system, application, network, configuration, process, or security control that could be used to cause an unwanted security outcome. Vulnerabilities can arise from insecure design, implementation errors, weak configurations, outdated software, missing security controls, or poor operational practices.

Understanding the relationship between vulnerabilities, threats, exploits, impacts, and risk is essential before performing vulnerability assessment or penetration testing. A vulnerability by itself does not always mean that an attack will occur; risk depends on factors such as the likelihood of exploitation and the potential impact.

\subsection{What is a Vulnerability?}
\subsubsection{Definition}
A vulnerability is a weakness, flaw, or condition in a system, application, network, configuration, process, or security control that can reduce the security of an asset. The weakness may exist because of insecure design, programming errors, incorrect configuration, outdated software, weak authentication, excessive privileges, or missing security controls. A vulnerability describes a condition that creates an opportunity for an unwanted security outcome; it does not by itself prove that an attack has occurred.

From a vulnerability-management perspective, weaknesses should be identified, validated, documented, prioritized, and addressed according to their technical and business context. The same weakness may require different levels of attention depending on the affected asset, exposure, available controls, and potential impact.

\subsubsection{Real-World Example}
A web application using an outdated third-party library with a publicly known security weakness contains a vulnerability. The weakness may remain present even if no attacker is currently attempting to use it.

\subsubsection{Cybersecurity Connection}
Vulnerability management focuses on discovering, validating, prioritizing, and reducing weaknesses before they can become significant security incidents.

\subsection{What is a Threat?}
\subsubsection{Definition}
A threat is a potential source, event, or circumstance that could cause harm to an asset. In cybersecurity, a threat can involve a malicious actor, an accidental action, a technical failure, or an environmental event. The important point is that a threat represents the possibility of harm, while a vulnerability represents a weakness that may contribute to that harm.

For example, an attacker attempting to compromise an exposed server is a threat. The threat becomes more significant when the target contains exploitable weaknesses, has valuable information, or is insufficiently protected by security controls. Threat awareness therefore helps security teams understand where vulnerabilities may be most relevant.

\subsubsection{Real-World Example}
An attacker attempting to compromise an exposed server is a threat to the organization operating that server.

\subsubsection{Cybersecurity Connection}
Threat information helps security teams understand which vulnerabilities and assets deserve greater attention.

\subsection{What is Risk?}
\subsubsection{Definition}
Risk represents the potential for loss or harm when a threat can take advantage of a vulnerability. Risk is therefore broader than the vulnerability itself. It considers factors such as the likelihood that an unwanted event could occur and the potential consequences if it does occur.

In practical cybersecurity work, risk is influenced by technical severity, exploitability, exposure, asset criticality, existing security controls, and business impact. This is why security teams do not always remediate every finding in exactly the same order. Risk-based prioritization helps direct limited resources toward weaknesses that could create the greatest overall concern.

\subsubsection{Real-World Example}
If a critical server has a remotely exploitable vulnerability and contains sensitive information, the likelihood and impact of exploitation may make the associated risk high.

\subsubsection{Cybersecurity Connection}
Risk-based prioritization allows security teams to address weaknesses according to their potential business and technical consequences rather than treating every finding equally.

\subsection{What is an Exploit?}
\subsubsection{Definition}
An exploit is a method, technique, or code that takes advantage of a vulnerability to produce an unintended security effect. An exploit may involve a specially crafted request, malicious input, a sequence of actions, or another technique designed to trigger the weakness.

The existence of an exploit does not automatically mean that every vulnerable system can be compromised in the same way. Exploitability depends on the affected software, configuration, privileges, network exposure, required conditions, and other security controls. In professional penetration testing, exploitation is performed only when it is explicitly authorized and appropriate for the defined scope.

\subsubsection{Real-World Example}
A crafted request that abuses a vulnerable application component to read unauthorized files is an example of exploitation.

\subsubsection{Cybersecurity Connection}
Penetration testing may validate exploitability in a controlled and explicitly authorized environment. Exploitation outside an approved scope is not permitted.

\subsection{What is Vulnerability Assessment?}
\subsubsection{Definition}
Vulnerability assessment is a systematic process of identifying, analyzing, prioritizing, and reporting security weaknesses in authorized systems. It may use automated scanners, configuration reviews, manual analysis, security advisories, and other assessment techniques to build an understanding of the security condition of the target.

The purpose is not simply to generate a long list of vulnerabilities. A useful assessment connects findings with affected assets, severity, evidence, potential impact, and remediation recommendations. The results can then support remediation planning and later verification through re-scanning or re-assessment.

\subsubsection{Expected Outcome}
The expected outcome is a structured list of findings with severity or risk information and remediation recommendations. Re-scanning can then be used to verify whether identified weaknesses have been addressed.

\fig{figure_01.png}{Relationship between vulnerability, threat, exploit, impact and risk}
\fig{figure_02.png}{Overview of the vulnerability assessment process}

\subsection{Vulnerability vs. Threat vs. Risk}
\begin{table}[H]\centering
\begin{tabularx}{\textwidth}{>{\bfseries}p{2.4cm}X X}
\toprule
Concept & Meaning & Simple Example\\
\midrule
Vulnerability & A weakness that can be abused or cause a security problem. & An outdated server component with a known security flaw.\\
Threat & A potential source of harm. & An attacker attempting to compromise an exposed server.\\
Risk & Potential loss or harm considering likelihood and impact. & A critical Internet-facing weakness on a sensitive server.\\
Exploit & A method used to take advantage of a vulnerability. & A crafted request that triggers a known vulnerable component.\\
\bottomrule
\end{tabularx}
\caption{Comparison of vulnerability, threat, risk and exploit}
\end{table}


\subsection{Quick Check}
\begin{enumerate}
\item What is the difference between a vulnerability and a threat?
\item Why does a vulnerability not automatically mean that a security incident will occur?
\item What factors commonly influence risk?
\item What is an exploit?
\item What is the purpose of vulnerability assessment?
\end{enumerate}

\subsection{Task 1 Summary}
Vulnerabilities are weaknesses, threats are potential sources of harm, exploits are methods used to abuse weaknesses, and risk represents potential loss or harm. These concepts form the foundation for vulnerability assessment and penetration testing.

\newpage
\section{Task 2: Common Security Vulnerabilities}

\subsection{SQL Injection (SQLi)}
\subsubsection{Definition}
SQL Injection occurs when untrusted input is incorporated into SQL statements without appropriate validation or parameterization. An attacker may manipulate the intended database query. In practice, the weakness occurs when application input crosses the boundary between ordinary data and database commands without adequate separation. Parameterized queries are important because they keep supplied values separate from the SQL statement structure.

\subsubsection{Effect on System}
A vulnerable application may process attacker-controlled input as part of a database command rather than as ordinary data.

\subsubsection{Possible Impact}
Potential impacts include unauthorized data access, data modification, authentication bypass, and in some situations broader compromise depending on the application's privileges and database configuration.

\subsubsection{Basic Prevention}
Use parameterized queries or prepared statements, validate input appropriately, apply least-privilege database permissions, and avoid constructing SQL statements through unsafe string concatenation.

\fig{figure_03.png}{SQL Injection attack flow through a vulnerable web application}

\refbasis{OWASP Top 10 – Injection}{https://owasp.org/Top10/A03_2021-Injection/}
\refstd{OWASP Web Security Testing Guide – SQL Injection Testing}{https://owasp.org/www-project-web-security-testing-guide/}
\video{Professor Messer – Web and Security Fundamentals}{https://www.youtube.com/@professormesser}

\subsection{Cross-Site Scripting (XSS)}
\subsubsection{Definition}
Cross-Site Scripting is a web application vulnerability in which attacker-controlled content is interpreted as script by a victim's browser. The underlying problem is that untrusted data reaches a browser-rendered context without appropriate handling. The exact defense depends on where the data is inserted, such as HTML, an attribute, JavaScript, CSS, or a URL context.

\subsubsection{Effect on System}
The malicious script can execute within the security context of the vulnerable web application when a victim accesses affected content.

\subsubsection{Possible Impact}
Depending on the application and browser context, impact can include manipulation of page content, unauthorized actions performed through the victim's session, or exposure of information accessible to client-side code.

\subsubsection{Basic Prevention}
Use context-appropriate output encoding, validate input where appropriate, apply safe templating practices, and use suitable browser security controls such as Content Security Policy.

\fig{figure_04.png}{Cross-Site Scripting (XSS) flow through a vulnerable web application}

\refbasis{OWASP Top 10 – Injection}{https://owasp.org/Top10/A03_2021-Injection/}
\refstd{OWASP Cross Site Scripting Prevention Cheat Sheet}{https://cheatsheetseries.owasp.org/cheatsheets/Cross_Site_Scripting_Prevention_Cheat_Sheet.html}
\video{Professor Messer – Web Security Concepts}{https://www.youtube.com/@professormesser}

\subsection{Broken Authentication}
\subsubsection{Definition}
Broken authentication refers to weaknesses in authentication or session-management controls that can allow unauthorized users to access accounts or protected functionality. Authentication weaknesses can occur at login, password recovery, session creation, session termination, or other identity-related functions. Because these controls establish and maintain user identity, weaknesses can directly affect access to protected resources.

\subsubsection{Effect on System}
Weak credential handling, poor session management, predictable tokens, or inadequate authentication controls can undermine account security.

\subsubsection{Possible Impact}
Possible impacts include account takeover, unauthorized access to protected resources, privilege misuse, and exposure of user information.

\subsubsection{Basic Prevention}
Use strong authentication controls, secure password storage, multi-factor authentication where appropriate, secure session handling, appropriate session expiration, and protection against credential attacks.

\fig{figure_05.png}{Broken authentication and its effect on access control}

\refbasis{OWASP Top 10 – Identification and Authentication Failures}{https://owasp.org/Top10/A07_2021-Identification_and_Authentication_Failures/}
\refstd{OWASP Authentication Cheat Sheet}{https://cheatsheetseries.owasp.org/cheatsheets/Authentication_Cheat_Sheet.html}
\video{Professor Messer – Authentication and Security Fundamentals}{https://www.youtube.com/@professormesser}

\subsection{Security Misconfiguration}
\subsubsection{Definition}
Security misconfiguration occurs when security settings are incorrectly implemented, left at insecure defaults, unnecessarily exposed, or inconsistently maintained. This category covers insecure settings in applications, operating systems, servers, databases, cloud resources, and network devices. It can occur when default settings remain enabled, unnecessary services are exposed, permissions are too broad, or security hardening is incomplete.

\subsubsection{Effect on System}
Examples include default credentials, unnecessary services, verbose error messages, exposed administrative interfaces, permissive permissions, and insecure server settings.

\subsubsection{Possible Impact}
Misconfiguration can increase the attack surface and may enable unauthorized access, information disclosure, or system compromise.

\subsubsection{Basic Prevention}
Use secure baseline configurations, remove unnecessary services, change default credentials, apply least privilege, disable unnecessary features, and review configurations regularly.

\fig{figure_06.png}{Security misconfiguration and its potential impact on application security}

\refbasis{OWASP Top 10 – Security Misconfiguration}{https://owasp.org/Top10/A05_2021-Security_Misconfiguration/}
\refstd{CIS Benchmarks}{https://www.cisecurity.org/cis-benchmarks}
\video{Professor Messer – Security Configuration Fundamentals}{https://www.youtube.com/@professormesser}

\subsection{Vulnerable and Outdated Components}
\subsubsection{Definition}
This vulnerability occurs when applications depend on libraries, frameworks, plugins, operating-system components, or other software with known or unaddressed security weaknesses. Modern applications depend on libraries, frameworks, plugins, operating-system packages, and other third-party components. If a component contains a known weakness and is not maintained, the application may inherit the security risk even when its own code has no obvious defect.

\subsubsection{Effect on System}
A vulnerable dependency can expose the application even when the application's own code is reasonably secure.

\subsubsection{Possible Impact}
Potential impact ranges from information disclosure and denial of service to unauthorized access or remote code execution, depending on the component and vulnerability.

\subsubsection{Basic Prevention}
Maintain an inventory of dependencies, monitor security advisories, apply supported updates, remove unnecessary components, and use dependency-management and software composition practices.

\fig{figure_07.png}{Vulnerable and outdated components leading to potential security impact}

\refbasis{OWASP Top 10 – Vulnerable and Outdated Components}{https://owasp.org/Top10/A06_2021-Vulnerable_and_Outdated_Components/}
\refstd{CISA Known Exploited Vulnerabilities Catalog}{https://www.cisa.gov/known-exploited-vulnerabilities-catalog}
\video{Professor Messer – Vulnerability and Patch Management Concepts}{https://www.youtube.com/@professormesser}

\subsection{Comparison of the Five Vulnerabilities}
\begin{table}[H]\centering
\small
\begin{tabularx}{\textwidth}{>{\bfseries}p{3.0cm}p{2.0cm}X X}
\toprule
Vulnerability & Typical Area & Possible Impact & Basic Prevention\\
\midrule
SQL Injection & Web/database & Unauthorized data access or modification & Parameterized queries; least privilege\\
XSS & Web application & Client-side script execution and session abuse & Output encoding; secure templating\\
Broken Authentication & Identity/session & Account takeover or unauthorized access & Strong authentication and session controls\\
Security Misconfiguration & Systems/apps & Increased attack surface and exposure & Secure baselines and configuration review\\
Vulnerable Components & Dependencies & Exploitation of known software weaknesses & Patch and maintain dependencies\\
\bottomrule
\end{tabularx}
\caption{Comparison of common security vulnerabilities}
\end{table}

\subsection{Quick Check}
\begin{enumerate}
\item Why is parameterized SQL preferred over string concatenation?
\item What is the role of output encoding in XSS prevention?
\item Name two controls that improve authentication security.
\item Give two examples of security misconfiguration.
\item Why should third-party components be maintained and updated?
\end{enumerate}


\subsection{Mini Practical – Vulnerability-to-Control Mapping}
\begin{table}[H]\centering
\small
\begin{tabularx}{\textwidth}{>{\bfseries}p{3.0cm}X X}
\toprule
Finding Type & Security Question & Example Control\\
\midrule
SQL Injection & Can input change database commands? & Parameterized queries\\
XSS & Can untrusted content execute in a browser? & Context-aware output encoding\\
Broken Authentication & Can an unauthorized user obtain valid access? & Strong authentication and session management\\
Security Misconfiguration & Are insecure settings increasing exposure? & Hardened baseline configuration\\
Outdated Component & Is a known vulnerable dependency deployed? & Patch and dependency management\\
\bottomrule
\end{tabularx}
\caption{Mapping common vulnerabilities to primary security controls}
\end{table}

\subsection{Task 2 Summary}
SQL Injection, XSS, broken authentication, security misconfiguration, and vulnerable components represent different classes of weaknesses. Preventive controls should address both technical implementation and ongoing security maintenance.

\newpage
\section{Task 3: Vulnerability Assessment Process}

\subsection{Introduction}
A vulnerability assessment follows a structured lifecycle so that findings are not only discovered but also analyzed, prioritized, remediated, and verified. The process should operate within an explicitly authorized scope.


\subsection{Step 1 – Identify}
\subsubsection{What Happens}
The assessor identifies assets, systems, applications, network ranges, services, and assessment scope.

\subsubsection{Why It Is Important}
An accurate asset inventory prevents important systems from being overlooked and ensures that testing stays within authorization boundaries.

\subsubsection{Expected Outcome}
A defined list of targets and an agreed assessment scope.


\subsection{Step 2 – Scan}
\subsubsection{What Happens}
Authorized scanning tools inspect the target for known vulnerabilities, exposed services, software versions, configuration weaknesses, and other indicators.

\subsubsection{Why It Is Important}
Automated scanning provides broad coverage and helps identify weaknesses that may require further validation.

\subsubsection{Expected Outcome}
A collection of potential findings requiring analysis.

\subsection{Step 3 – Analyze}
\subsubsection{What Happens}
Findings are reviewed for accuracy, context, affected assets, severity, and evidence.

\subsubsection{Why It Is Important}
Automated scanners can produce false positives or findings whose importance depends on the environment.

\subsubsection{Expected Outcome}
Validated and contextualized findings.

\subsection{Step 4 – Prioritize}
\subsubsection{What Happens}
Findings are ranked using severity, likelihood, exploitability, exposure, asset importance, and potential impact.

\subsubsection{Why It Is Important}
Prioritization allows security teams to address the most important risks first.

\subsubsection{Expected Outcome}
A risk-ranked remediation list.

\subsection{Step 5 – Remediate}
\subsubsection{What Happens}
The organization applies patches, configuration changes, access controls, compensating controls, or other corrective measures.

\subsubsection{Why It Is Important}
Identification alone does not improve security unless weaknesses are reduced or eliminated.

\subsubsection{Expected Outcome}
The identified weakness is corrected or its risk is reduced to an acceptable level.


\subsection{Step 6 – Re-scan and Verify}
\subsubsection{What Happens}
The affected system is assessed again to determine whether the vulnerability has been remediated.

\subsubsection{Why It Is Important}
Verification provides evidence that remediation was effective.

\subsubsection{Expected Outcome}
A verified remediation status and an updated security assessment.

\fig{figure_08.png}{Vulnerability assessment lifecycle from identification to verification}

\subsection{Quick Check}
\begin{enumerate}
\item Why is asset identification important?
\item What is the purpose of vulnerability scanning?
\item Why should findings be analyzed before remediation?
\item What factors can influence prioritization?
\item Why is re-scanning necessary?
\end{enumerate}

\subsection{Task 3 Summary}
The vulnerability assessment process follows Identify → Scan → Analyze → Prioritize → Remediate → Re-scan and Verify. This lifecycle turns scanner output into an actionable and verifiable security-improvement process.

\newpage
\section{Task 4: Introduction to Penetration Testing}


\subsection{Scope, Rules of Engagement and Authorization}
A penetration test should begin with a clear agreement between the tester and the system owner.

\subsubsection{Scope}
Scope identifies the systems, applications, addresses, accounts, environments, and testing windows that are included.

\subsubsection{Out-of-Scope Systems}
Explicit exclusions prevent accidental testing of systems that the tester is not authorized to assess.

\subsubsection{Testing Constraints}
The engagement may restrict techniques that could cause service interruption, data modification, destructive changes, or uncontrolled propagation.

\subsection{Why Written Permission Is Important}
Security testing creates traffic and may involve techniques that resemble malicious activity. Written authorization provides evidence that the activity is approved and establishes the boundaries of the assessment.

\subsubsection{Authorization Details}
A written authorization should identify the responsible organization or owner, the approved targets, assessment dates, tester or team, and relevant conditions.

\subsubsection{Emergency Contact}
The tester should know whom to contact if the assessment causes an unexpected service issue or discovers a severe security condition requiring immediate attention.

\subsection{VA and PT – Practical Difference}
A vulnerability assessment can identify many potential weaknesses across a broad environment. A penetration test generally selects relevant weaknesses or attack paths and validates their practical security significance under controlled conditions.

\subsubsection{Coverage}
Vulnerability assessment generally emphasizes breadth and identification. Penetration testing emphasizes validation, attack-path reasoning, and impact.

\subsubsection{Output}
A vulnerability assessment commonly produces prioritized findings and remediation recommendations. A penetration test may additionally document validated attack paths, evidence of exploitation, security-control observations, and business impact.

\subsection{Ethical and Professional Considerations}
A security tester must protect client information, minimize unnecessary impact, stay within scope, document actions accurately, and report findings responsibly. Professional testing is intended to improve security rather than create uncontrolled damage.

\subsection{What is Penetration Testing?}
\subsubsection{Definition}
Penetration testing is an authorized security assessment in which a tester evaluates whether security weaknesses can be practically abused within a defined scope. It is more than simply listing vulnerabilities: it focuses on validating security controls and understanding potential impact.

\subsubsection{Why Penetration Testing Is Performed}
Penetration testing can help organizations validate defenses, identify exploitable weaknesses, assess attack paths, test detection and response capabilities, and produce evidence-based remediation recommendations.

\subsection{Vulnerability Assessment vs. Penetration Testing}
Vulnerability assessment primarily focuses on discovering and prioritizing weaknesses. Penetration testing goes further by validating selected weaknesses through controlled testing within the approved scope.

\fig{figure_14.png}{Comparison of vulnerability assessment and penetration testing}

\subsection{Authorized Security Testing}
\subsubsection{Written Permission}
Written permission establishes that the tester is authorized to assess the specified systems. It also helps define scope, time windows, permitted techniques, contacts, and handling requirements.

\subsubsection{Why Authorization Matters}
Scanning, testing, or exploiting systems without permission can disrupt services, expose data, violate policy, or create legal consequences. Professional security testing therefore begins with authorization and a clearly defined scope.


\subsection{Penetration Testing Lifecycle}
\subsubsection{Planning}
Define objectives, scope, authorization, constraints, contacts, and reporting requirements.

\subsubsection{Reconnaissance}
Collect permitted information about the target environment to understand the attack surface.

\subsubsection{Scanning}
Use appropriate authorized techniques to identify services, technologies, and potential weaknesses.

\subsubsection{Vulnerability Analysis}
Analyze discovered information and determine which weaknesses warrant controlled validation.

\subsubsection{Security Testing}
Perform controlled testing against approved targets while minimizing disruption and avoiding unnecessary impact.

\subsubsection{Reporting}
Document findings, evidence, risk, business/technical impact, and remediation recommendations.

\fig{figure_09.png}{Penetration testing lifecycle from planning to security reporting}

\subsection{Authorized Penetration Testing Environment}
A professional assessment should use an explicitly authorized target such as an organization's approved system, the tester's own computer, an intentionally vulnerable laboratory, or an approved training platform.

\fig{figure_10.png}{Authorized penetration testing workflow within a defined security scope}

\refbasis{NIST SP 800-115 – Technical Guide to Information Security Testing and Assessment}{https://csrc.nist.gov/publications/detail/sp/800-115/final}
\refstd{NIST SP 800-115}{https://csrc.nist.gov/publications/detail/sp/800-115/final}
\video{Professor Messer – Penetration Testing and Security Assessment Concepts}{https://www.youtube.com/@professormesser}

\subsection{Quick Check}
\begin{enumerate}
\item What is the main difference between vulnerability assessment and penetration testing?
\item Why is written permission important?
\item What is defined during the planning phase?
\item Why should penetration testing remain within a defined scope?
\item What should be included in the final report?
\end{enumerate}

\subsection{Task 4 Summary}
Penetration testing is an authorized, controlled security assessment. A professional lifecycle begins with planning and authorization, progresses through reconnaissance, scanning, vulnerability analysis and controlled testing, and concludes with evidence-based reporting.

\newpage
\section{Task 5: Safe Practical Activity}

\subsection{Introduction}
For this practical activity, Nessus Essentials was used as the primary vulnerability assessment tool in a controlled laboratory environment. The target was an intentionally vulnerable Metasploitable system. The activity was performed for security learning and assessment purposes within the lab environment.

\subsection{Tool Used – Nessus Essentials}
Nessus is a vulnerability assessment platform used to identify security weaknesses in systems and applications. In this practical, Nessus was configured as a local scanner and used to perform a Basic Network Scan against the laboratory target.

\fig{nessus_setup.png}{Nessus setup screen used to configure the Nessus deployment}
\fig{nessus_events.png}{Nessus Events page showing successful component and plugin downloads}


\subsection{Practical Objective}
The objective of this activity was to perform a controlled vulnerability assessment against an intentionally vulnerable Metasploitable laboratory target using Nessus Essentials. The exercise focused on identifying vulnerabilities, reviewing scanner evidence, understanding severity information, and documenting remediation recommendations.

\subsection{Why Nessus Was Selected}
Nessus is appropriate for this practical because it is specifically designed for vulnerability assessment. It provides structured findings, severity information, plugin identifiers, evidence, risk metrics, and remediation guidance. This makes it suitable for demonstrating the complete assessment workflow required by the task.

\subsubsection{Assessment Workflow}
The practical followed the sequence:
\begin{enumerate}
\item Prepare and verify the Nessus installation.
\item Confirm the authorized laboratory target.
\item Create the vulnerability assessment scan.
\item Execute the scan using the selected policy.
\item Review the completed findings.
\item Examine representative Critical findings.
\item Record evidence, risk and remediation.
\end{enumerate}

\fig{figure_11.png}{Nessus-based vulnerability assessment workflow}

\subsection{Nessus Assessment Observations}
The Nessus interface provided a completed scan report with the number of vulnerabilities and their severity distribution. Individual plugin pages supplied additional information such as description, solution, plugin ID, CVSS, VPR, EPSS, affected host and observed output.

\subsubsection{Severity Distribution}
The results interface displayed Critical, High, Medium, Low, and Informational categories. These categories provide a convenient first-level method of organizing findings for review.

\subsubsection{Plugin-Level Evidence}
The plugin detail view provides information that is useful for a security finding report. This includes the plugin identifier, severity, description, solution, risk information, and output observed during the scan.

\subsection{Observed Vulnerability Results}
The completed scan reported **72 vulnerabilities**. The Nessus results interface classified findings by severity and displayed additional metrics such as CVSS, VPR and EPSS where available.

\fig{nessus_results.png}{Completed Nessus vulnerability assessment showing the observed findings and severity distribution}

\subsection{Finding 1 – UnrealIRCd Backdoor Detection}
\subsubsection{Finding}
Nessus identified **UnrealIRCd Backdoor Detection** as a Critical finding, Plugin ID **46882**.

\subsubsection{Risk Level}
**Critical**. The Nessus result displayed a CVSS v3.0 base score of **10.0** and a VPR value of **7.7** for this finding.

\subsubsection{Potential Impact}
The plugin description states that the affected remote IRC service contains a backdoor that could allow arbitrary code execution on the affected host. In the observed scan output, the service was identified on port **6667/tcp**.

\subsubsection{Recommendation}
The Nessus solution recommends re-downloading the affected software, verifying its integrity using published checksums, and reinstalling it. In a real environment, affected software should be removed or upgraded according to the vendor's supported remediation guidance.

\fig{nessus_unrealircd.png}{Nessus detail view for the Critical UnrealIRCd Backdoor Detection finding}

\subsection{Finding 2 – Apache Tomcat AJP Connector Request Injection}
\subsubsection{Finding}
Nessus identified **Apache Tomcat AJP Connector Request Injection (Ghostcat)** as a Critical finding, Plugin ID **134862**.

\subsubsection{Risk Level}
**Critical**. The displayed CVSS v3.0 base score was **9.8**, with a VPR of **7.9** and EPSS of **0.9927** in the observed result.

\subsubsection{Potential Impact}
The finding description indicates that a remote unauthenticated attacker could potentially read web application files from a vulnerable AJP connector. In conditions where file upload is possible, the weakness can have more serious consequences.

\subsubsection{Recommendation}
The Nessus solution recommends configuring the AJP connector to require authorization and/or upgrading the affected Tomcat version to a supported version that addresses the vulnerability.

\fig{nessus_ghostcat.png}{Nessus detail view for the Critical Apache Tomcat AJP Connector Request Injection finding}


\subsection{Security and Safety Rule}
Scanning, testing, exploiting, or attacking systems without permission is not part of this practical activity. The exercise must remain restricted to the authorized laboratory environment.

\refbasis{Tenable Nessus – Vulnerability Assessment Platform}{https://www.tenable.com/products/nessus}
\refstd{Tenable Nessus Essentials}{https://www.tenable.com/products/nessus/nessus-essentials}
\refstd{OWASP Vulnerable Web Applications Directory – Training Environments}{https://vwad.owasp.org/}
\video{Tenable – Nessus Vulnerability Management and Assessment Resources}{https://www.youtube.com/@TenableSecurity}

\subsection{Quick Check}
\begin{enumerate}
\item What was the purpose of using Nessus in this practical?
\item What target was assessed?
\item How many vulnerabilities were reported by the completed scan?
\item What severity was assigned to the UnrealIRCd Backdoor Detection finding?
\item Why must vulnerability scanning be restricted to authorized targets?
\end{enumerate}

\subsection{Task 5 Summary}
The practical demonstrated a controlled vulnerability assessment using Nessus Essentials against an intentionally vulnerable Metasploitable system. The completed scan produced 72 findings and provided severity, scoring, evidence, and remediation information for individual vulnerabilities.

\newpage
\section{Task 6: Security Finding Report}

\subsection{Introduction}
A security finding report converts technical assessment results into information that can be understood and acted upon by system administrators, developers, security teams, and management. Each finding should clearly state the weakness, risk level, potential impact, and recommended remediation.

\fig{figure_13.png}{Structure of a security finding report}

\subsection{Finding 1 – UnrealIRCd Backdoor Detection}
\begin{table}[H]\centering
\begin{tabularx}{\textwidth}{>{\bfseries}p{3.0cm}X}
\toprule
Finding & UnrealIRCd Backdoor Detection (Plugin 46882)\\
Risk Level & Critical\\
Possible Impact & Potential arbitrary code execution on the affected host through the vulnerable backdoored service.\\
Recommendation & Re-download and verify the software using trusted checksums, reinstall a clean supported version, and remove the affected backdoored component.\\
Evidence & Nessus identified the finding on the laboratory target and displayed the affected service on TCP port 6667.\\
\bottomrule
\end{tabularx}
\caption{Security finding report – UnrealIRCd Backdoor Detection}
\end{table}

\subsection{Finding 2 – Apache Tomcat AJP Connector Request Injection}
\begin{table}[H]\centering
\begin{tabularx}{\textwidth}{>{\bfseries}p{3.0cm}X}
\toprule
Finding & Apache Tomcat AJP Connector Request Injection (Ghostcat), Plugin 134862\\
Risk Level & Critical\\
Possible Impact & Potential unauthorized reading of web application files through a vulnerable AJP connector; impact can be greater where vulnerable configurations permit file upload.\\
Recommendation & Require authorization for the AJP connector and upgrade the affected Tomcat deployment to a supported version addressing the vulnerability.\\
Evidence & Nessus identified the finding and provided plugin-level description, scoring, output, and remediation information.\\
\bottomrule
\end{tabularx}
\caption{Security finding report – Apache Tomcat AJP Connector Request Injection}
\end{table}


\subsection{How to Read a Nessus Finding}
A Nessus finding can contain several technical fields. The finding name identifies the weakness, the plugin ID provides a unique scanner reference, severity provides an initial classification, and scoring fields such as CVSS, VPR, or EPSS provide additional risk context.

\subsubsection{CVSS}
CVSS provides a standardized way of expressing the technical severity of a vulnerability using characteristics such as attack vector, complexity, privileges, user interaction, scope, and impact.

\subsubsection{VPR}
Vulnerability Priority Rating is a Tenable risk-prioritization metric that can help determine which vulnerabilities deserve attention.

\subsubsection{EPSS}
Exploit Prediction Scoring System information estimates the likelihood of exploitation in the wild and can provide additional context for prioritization.

\subsection{Finding Documentation Quality}
Each reported finding should be understandable without requiring the reader to inspect the entire scanner interface.

\subsubsection{Finding}
State exactly what weakness was detected and identify the relevant component.

\subsubsection{Risk Level}
Record the severity observed in the assessment and explain any important contextual factors.

\subsubsection{Potential Impact}
Describe what an attacker could potentially achieve if the weakness were successfully abused.

\subsubsection{Recommendation}
Provide a practical corrective action such as upgrading software, changing configuration, disabling an unnecessary service, strengthening authentication, or applying an appropriate security control.

\subsection{Remediation Priority}
Critical findings generally receive immediate attention, but final priority should still consider exposure and asset context. High and Medium findings should be tracked through the organization's remediation process, while Low and Informational findings may still provide useful hardening opportunities.

\subsection{Vulnerability Finding and Risk Prioritization}
The assessment process should validate findings, evaluate likelihood and impact, assign a priority, and select an appropriate remediation action.

\fig{figure_12.png}{Vulnerability finding analysis and risk prioritization workflow}

\subsection{Security Risk Rating}
Risk prioritization should consider both the likelihood that a weakness can be exploited and the potential impact on the affected system or organization. Severity ratings from scanning tools provide useful technical context but should be interpreted alongside asset criticality and exposure.

\fig{figure_16.png}{Security risk evaluation based on likelihood and potential impact}

\subsection{Task 6 Summary}
A useful security finding report connects each weakness to its risk, possible impact, evidence, and remediation. The practical findings from Nessus demonstrate how scanner output can be converted into actionable security recommendations.

\subsection{Final Week 4 Summary}
This Week 4 assignment provided a structured introduction to Vulnerability Assessment and Penetration Testing. The report covered the concepts of vulnerabilities, threats, risks, exploits, common security weaknesses, the vulnerability assessment lifecycle, penetration testing methodology, authorization requirements, and safe practical assessment.

The practical activity demonstrated the use of Nessus Essentials against an intentionally vulnerable Metasploitable laboratory system. The completed scan identified 72 vulnerabilities, including Critical findings such as UnrealIRCd Backdoor Detection and Apache Tomcat AJP Connector Request Injection. These results were documented with risk, impact, evidence, and remediation recommendations.

The overall learning outcome is that professional cybersecurity assessment is not simply about finding vulnerabilities. It requires authorized scope, systematic assessment, careful analysis, risk-based prioritization, safe testing, evidence collection, remediation, and verification.


\end{document}
